\section{A universal monomial obstruction for nonsplit sixfold sources}
\label{sec:monomial-sixfold}

Here ``monomial'' refers to induction of a complex representation from a
one-dimensional character. It does not require a reflection group or
reflection inertia. We allow the values of the inducing character to
generate a cyclotomic field strictly larger than the quadratic character
field. The conclusion concerns primitive factors of the rational group
algebra; it does not concern an unrelated quadratic action imposed on a
factor whose character field is real.

\subsection*{D-MON-01. The cyclic factor, including its Schur index}

Let $G$ act faithfully on a connected smooth projective curve $C$, and let
$\rho$ be an irreducible complex representation with character field an
imaginary quadratic field $K$. Suppose
\[
             \rho=\operatorname{Ind}_H^G\chi,
             \qquad d=\operatorname{ord}(\chi),
             \qquad L=\mathbf Q(\zeta_d).
\]
Let $B$ be the abelian factor selected by a primitive idempotent in the
corresponding simple component of $\mathbf Q[G]$. Assume that $B$ is a
sixfold with $K$-signature $(3,3)$, with the polarization induced from
$J(C)$. The character formula gives $K\subset L$, and $d>2$ because a
real-valued linear character would induce a rational-valued character.

Write that simple component as $M_k(D)$, where $D$ is a central division
algebra over $K$, and let $s^2=\dim_KD$ be its Schur-index square. The
corner algebra of a primitive idempotent acts faithfully on $H^1(B,\mathbf Q)$.
Consequently this cohomology is a module over $D$, or its opposite according
to the cohomological convention, and
\[
                     s^2\mid\dim_KH^1(B,\mathbf Q)=6.
\]
The only possible positive integer $s$ is one. Thus $\rho$ has a
$K$-realization $V$, its component is $M_{\dim\rho}(K)$, and the
multiplicity Hodge structure
\[
     M=\operatorname{Hom}_{\mathbf Q[G]}(V_{\mathbf Q},H^1(C,\mathbf Q))
\]
is the rational Hodge structure of $B$. Here $V_{\mathbf Q}$ means
restriction of scalars from $K$; its commuting algebra is $K$.

Put $J=\ker\chi$, $Y=C/J$ and $X=C/H$. Then $Y\to X$ is a connected
cyclic cover of degree $d$. Let $P_d(Y/X)$ denote the primitive
cyclotomic factor of $J(Y)$, namely the image of the rational idempotent
for $\mathbf Q(\zeta_d)$ in $\mathbf Q[H/J]$. There is an isogeny
\begin{equation}\label{eq:mon-power}
           P_d(Y/X)\ \sim\ B^{\ell},\qquad \ell=[L:K].
\end{equation}
Here is the rational representation argument, including the multiplicity.
Regard $L$ as the rational $H$-module afforded by $\chi$. Over $\mathbf C$,
the induction of this module is the sum of $\operatorname{Ind}_H^G\chi^\sigma$
over all embeddings $\sigma:L\hookrightarrow\mathbf C$. Every summand is
$\rho$ or $\bar\rho$, according to the restriction of $\sigma$ to $K$,
and each occurs $[L:K]$ times. Since the Schur index is one, semisimplicity
over $\mathbf Q$ therefore gives
\[
             \operatorname{Ind}_H^G L\cong V_{\mathbf Q}^{\oplus\ell}.
\]
Frobenius reciprocity identifies its Hom into $H^1(C,\mathbf Q)$ with
$\operatorname{Hom}_H(L,H^1(C,\mathbf Q))$. The latter is precisely
the primitive cyclotomic Hodge structure of $Y$: first take $J$-invariants,
then the primitive $\mathbf Q[H/J]$ summand. This proves
\eqref{eq:mon-power}. These constructions use fixed rational group-algebra
maps, so in a family they give an isomorphism of rational variations and,
after clearing fixed denominators, an isogeny of abelian families.

Write $h=g(X)$ and let $R$ be the number of branch points of $Y\to X$.
The faithful character $\chi$ has no invariants at any genuine branch
point. Its parabolic Euler characteristic and its $(3,3)$ multiplicities
give
\begin{equation}\label{eq:mon-dimension}
                  6=2h-2+R,\qquad R=8-2h.
\end{equation}
One can also obtain this identity by adding the two character
Chevalley--Weil formulas: each nonzero local exponent and its conjugate
contribute one. In particular $0\le h\le4$.

For fixed cyclic branch exponents, the moduli of these cyclic covers are
finite over the moduli of the $R$-pointed base curve. Indeed the cover is
specified by a $d$th root of the appropriate branch line bundle, and roots,
when they exist, form a finite torsor under $J(X)[d]$. Using
\eqref{eq:mon-dimension}, their dimensions for $h=0,1,2,3,4$ are at most
\[
                              5,6,7,8,9,
\]
respectively. The genus-zero and genus-one formulas here use $R=8$ and
$R=6$, so no unstable pointed-curve exception occurs.

The period differential of $B$ has the same kernel as that of $B^\ell$:
on the Hodge bundles the latter is the direct sum of $\ell$ copies of the
former. Isogeny \eqref{eq:mon-power} therefore bounds its rank by the
rank of the forgetful map to cyclic-cover moduli. This is a statement
about actual variations, not a sufficiency claim based on the source
dimension. If the $B$-period map has rank nine, necessarily
\begin{equation}\label{eq:mon-equality}
                  h=4,\qquad R=0,
\end{equation}
and the forgetful image is dominant in a component of the moduli of
connected etale cyclic degree-$d$ covers of genus-four curves. After
passing to finite level structures these are algebraic varieties; a
dominant constructible image contains a nonempty Zariski open set.
No assumption about simple branching of $C\to C/G$ is used here.

\subsection*{D-MON-02. The full endomorphism algebra of a general cyclic Prym}

For a connected etale cyclic degree-$d$ cover of a genus-four curve, the
primitive integral homology, modulo the harmless choice of a finite-index
lattice, is a rank-six module over $R_d=\mathbf Z[\zeta_d]$.
Looijenga's Proposition~2.2 identifies its equivariant intersection form
with the standard skew-Hermitian hyperbolic form
\[
               \Omega=\begin{pmatrix}0&I_3\\-I_3&0\end{pmatrix}.
\]
His Theorem~2.4 says that the group of lifts of mapping classes fixing
the cyclic quotient has image
\[
 U_6^\#(R_d)=
 \{g\in\operatorname{GL}_6(R_d):g^*\Omega g=\Omega,
                       \ \det g\in\langle\zeta_d^2\rangle\}.
\]
Both statements occur on page~5 of the author's version of
\cite{IM:Looijenga1997}; the hyperbolic basis is proved on pages~8--9,
and the genus-at-least-three monodromy proof ends on page~12.
We use the published theorem, and do not claim a new proof of its full
mapping-class surjectivity. A finite level cover has monodromy of finite
index in this image, which is all that is needed below.

We first prove the precise commutant assertion. For every finite-index
subgroup $\Gamma\subset U_6^\#(R_d)$,
\begin{equation}\label{eq:mon-commutant}
               \operatorname{End}_{\mathbf Q,\Gamma}(L^6)=L.
\end{equation}
Indeed $U_6^\#(R_d)$ contains $\operatorname{Sp}_6(\mathbf Z)$.
For every integral symplectic transvection $1+N$, some positive power
belongs to $\Gamma$. As $N^2=0$, commuting with that power is equivalent
over $\mathbf Q$ to commuting with $N$. Thus a commuting endomorphism
commutes with all integral symplectic transvections. The transvections
in the coordinate vectors give the elementary matrices
$E_{i,i+3},E_{i+3,i}$, whose products give the coordinate projectors.
Transvections in sums of coordinate vectors identify the resulting
coordinate endomorphisms. Hence the commutant on $L^6$, viewed as a
rational vector space, consists of operators
$\operatorname{id}_{\mathbf Q^6}\otimes f$ with
$f\in\operatorname{End}_{\mathbf Q}(L)$.

Now let $B$ be the Hermitian three-by-three matrix with
$B_{12}=\zeta_d$, $B_{21}=\bar\zeta_d$ and all other entries zero.
The matrix
\[
                          u_B=\begin{pmatrix}I_3&B\\0&I_3\end{pmatrix}
\]
belongs to $U_6^\#(R_d)$ and is unipotent with square-zero logarithm.
Some positive power belongs to $\Gamma$. Commutation with that power
forces $f(\zeta_d x)=\zeta_d f(x)$ for every $x\in L$. Thus $f$ is
$L$-linear and is multiplication by $f(1)$. This proves
\eqref{eq:mon-commutant}, without invoking an unproved assertion about
Zariski closure or independence of Galois embeddings.

It follows that the primitive abelian variety has full rational
endomorphism algebra $L$ outside a countable union of proper analytic
loci. We give the argument because a statement only about endomorphisms
commuting with the cyclic group would not suffice for our application.
On the connected universal cover of a finite-level moduli component,
trivialize rational homology. For each of its countably many rational
endomorphisms, the condition of preserving the Hodge subbundle is closed
analytic. Outside the union of those loci which are proper, the Hodge
endomorphisms are exactly those preserving the Hodge structure everywhere
on this universal cover. This rational algebra $E$ is stable under
monodromy conjugation, and its intersection with integral endomorphisms
is a lattice. The polarization gives the positive definite Rosati norm
\[
                         f\longmapsto\operatorname{Tr}(f^\dagger f)
\]
on this lattice. Monodromy conjugation preserves the lattice and this
norm, so its action on $E$ is finite. A finite-index monodromy subgroup
therefore centralizes $E$. Equation \eqref{eq:mon-commutant} gives
$E\subset L$; the reverse inclusion is the cyclic deck action. Hence
$E=L$. The complement of the countably many proper analytic loci meets
every nonempty analytic open set, by the Baire theorem. This is the only
genericity consequence needed below; an algebraicity theorem for Hodge
loci is not needed.

\subsection*{D-MON-03. The universal monomial sixfold barrier}

Under the hypotheses of D-MON-01, a rank-nine period map forces
\begin{equation}\label{eq:mon-final}
                 L=K,\qquad d\in\{3,4,6\},
                 \qquad\text{the polarization Hermitian space is split.}
\end{equation}
In particular there is no dominant nonsplit Weil-sixfold curve-cover
source from any such monomial representation. For a quadratic field
other than $\mathbf Q(i)$ and $\mathbf Q(\sqrt{-3})$, there is no
dominant source of this kind even in the split component.

To prove this, rank nine first gives \eqref{eq:mon-equality} and a
nonempty open image in the etale cyclic genus-four moduli. Choose a
point of this open image at which D-MON-02 gives
$\operatorname{End}^0(P_d)=L$. If $\ell=[L:K]>1$, the isogeny
$P_d\sim B^\ell$ gives an embedding
\[
                       M_\ell(K)\hookrightarrow\operatorname{End}^0(P_d).
\]
The left side is noncommutative and the right side is a field, a
contradiction. Thus $L=K$, so the cyclotomic degree is two and
$d=3,4,6$.

For completeness the final splitness concerns the actual induced
polarization. When $L=K$, the rational element
\[
 p_\chi=\frac1{|H|}\sum_{h\in H}
                   \bigl(\chi(h)^{-1}+\bar\chi(h)^{-1}\bigr)h
                   \in\mathbf Q[G]
\]
is a Rosati-fixed primitive idempotent of the $\rho$ component.
Frobenius reciprocity gives multiplicity one of $\chi$ in
$\rho|_H$, and induction shows that no other complex irreducible
representation contains $\chi$. Thus this idempotent selects exactly
one $K$-copy. Its abelian image is the pullback of $P_d(Y/X)$ along
$C\to Y$; the polarization pulls back to $|J|$ times the restriction
of the Jacobian polarization on $P_d$.

Looijenga's equivariant hyperbolic basis, or the cut-surface construction
in its proof, says that the latter skew-Hermitian space is the sum of
three hyperbolic planes. Multiplication by a nonzero imaginary scalar
converts it to a Hermitian form with the same half-dimensional isotropic
subspace. The actual polarization Hermitian form is therefore split.
Any other Rosati-fixed primitive factor in the same component has the
same multiplicity Hermitian form multiplied by the positive rational
norm of a line in the positive definite representation space. This
preserves splitness. Equivalently, restrict the tensor-product
polarization on the isotypic decomposition to either primitive line.
The raw determinant class is $[-1]$ in dimension six, or the trivial
class in the signed-discriminant convention.

This conclusion allows arbitrary initial inertia orders, arbitrary base
genus for $C/G$, arbitrary Nielsen class, and every monomial induction
whose character field is $K$. Its necessary equality condition was
deduced from the actual cyclic variation. No forgetful submersion for
fixed higher ramification was assumed. The remaining finite-group
nonsplit-source gate must use a nonmonomial irreducible representation,
or must change the specified group-algebra/character-field setup.
